\section{The numen: finite digits, infinite digits, and place-dependent limits}
\label{sec:numen}

The numen is the translation cocycle of the affine word representation.  Its
extension to infinite digits is a second construction, and the topology in
which that extension converges is part of the data.  Keeping these two facts
separate removes several ambiguities present when the word ``continuation''
is used without a codomain place.

\subsection{Finite base-\texorpdfstring{$p$}{p} words}

For $\mathbf j=(j_0,\ldots,j_{m-1})\in\Sigma_p^*$ define
\begin{equation}
\label{eq:digsum}
 \Dig(\mathbf j)=\sum_{a=0}^{m-1}j_ap^a.
\end{equation}
The least significant digit occurs first.  The map $\Dig:\Sigma_p^*\to
\N_0$ is surjective and is not injective: appending zeroes on the right does
not change its value.  These are the only identifications.

\begin{lemma}[finite-word quotient]
\label{lem:word-quotient}
If $\Dig(\mathbf i)=\Dig(\mathbf j)$, then one of $\mathbf i,\mathbf j$ is
obtained from the other by appending finitely many zeroes, after deleting
zeroes already appended to both.  Consequently
$\Sigma_p^*/\!\sim\;\cong\N_0$, where
$\mathbf i\sim\mathbf j$ means equality of digit sums.
\end{lemma}

\begin{proof}
Every nonzero integer has a unique finite base-$p$ expansion with nonzero
last digit.  The zero integer is represented by the empty word and by words
consisting only of zeroes.  The claim follows.
\end{proof}

\begin{proposition}[descent of the word numen]
\label{prop:numen-descent}
If $H$ is centered, then $X_H(\mathbf j)$ depends only on
$\Dig(\mathbf j)$.  Hence there is a well-defined function
\[
 X_H:\N_0\longrightarrow K,
 \qquad X_H(\Dig(\mathbf j))=X_H(\mathbf j).
\]
\end{proposition}

\begin{proof}
Centeredness gives $H_0(0)=0$.  Therefore
\[
 X_H(\mathbf j\wedge(0))
 =H_{\mathbf j}(H_0(0))=H_{\mathbf j}(0)=X_H(\mathbf j).
\]
Iterate this identity and use \cref{lem:word-quotient}.
\end{proof}

The multiplier requires a separate convention: unlike the centered numen, it
does not descend through the finite-word quotient.

\begin{definition}[canonical integer multiplier]
\label{def:canonical-integer-multiplier}
For $n>0$, let
\[
 \operatorname{dig}_p(n)=(d_0(n),\ldots,d_{L-1}(n))
\]
be the unique shortest base-$p$ word representing $n$, so that
$d_{L-1}(n)\ne0$, and put $\operatorname{dig}_p(0)=\varnothing$.  Define
\begin{equation}
\label{eq:canonical-integer-multiplier}
 M_H(n):=M_H\bigl(\operatorname{dig}_p(n)\bigr),
 \qquad M_H(0)=I.
\end{equation}
Thus $M_H(n)$ always refers to the canonical word, not to an arbitrary word
with digit sum $n$.
\end{definition}

Indeed, the concatenation identity gives
\[
 M_H\bigl(\mathbf j\wedge(0)^{\wedge k}\bigr)
 =M_H(\mathbf j)r_0^k.
\]
Consequently equivalent representatives obtained by appending zeroes need
not have the same multiplier.  The word multiplier descends through all such
representatives if and only if $r_0=I$.

For a noncentered proper Hydra, the natural digit-recursive function still
exists, but it is not the translation term $H_{\mathbf j}(0)$ alone.  Its
terminal value must be the fixed point of the zero branch.

\begin{theorem}[digit-recursive characterization]
\label{thm:digit-recursive-characterization}
Let $(r_j,c_j)_{0\leq j<p}$ be affine branches on $K$.  Functions
$f:\N_0\to K$ satisfying
\begin{equation}
\label{eq:numen-FE-N}
 f(pn+j)=r_jf(n)+c_j
 \qquad(n\in\N_0,\ 0\leq j<p)
\end{equation}
are in bijection with solutions $a\in K$ of
\begin{equation}
\label{eq:zero-fixed-equation}
 (I-r_0)a=c_0.
\end{equation}
The bijection sends $f$ to $f(0)=a$.  If $H$ is proper, the solution is
unique and $a=(I-r_0)^{-1}c_0$.  If $H$ is centered and proper, $a=0$ and
$f=X_H$.
\end{theorem}

\begin{proof}
Putting $n=j=0$ in \cref{eq:numen-FE-N} gives
\cref{eq:zero-fixed-equation}.  Conversely, choose a solution $a$.  Set
$f(0)=a$.  Every positive $n$ has a unique decomposition $n=pm+j$ with
$0\leq j<p$ and $m<n$, so \cref{eq:numen-FE-N} recursively defines one and
only one value $f(n)$.  The recursion automatically satisfies every
displayed equation because the decomposition is unique.

If $\mathbf d=\operatorname{dig}_p(n)=(d_0,\ldots,d_{L-1})$, repeated
substitution gives
\begin{equation}
\label{eq:finite-general-solution}
 f(n)=H_{d_0}\circ\cdots\circ H_{d_{L-1}}(a)
     =M_H(\mathbf d)a+X_H(\mathbf d).
\end{equation}
For centered $H$, appending zeroes leaves the value unchanged and this is
exactly the descended word numen of \cref{prop:numen-descent}.
\end{proof}

Combining \cref{eq:X-word} with the canonical digit word gives the explicit
finite formula
\begin{equation}
\label{eq:numen-finite-digit-formula}
 X_H(n)=\sum_{a=0}^{L-1}
 \bigl(r_{d_0}\cdots r_{d_{a-1}}\bigr)c_{d_a},
 \qquad \operatorname{dig}_p(n)=(d_0,\ldots,d_{L-1}),
\end{equation}
for a centered Hydra, with the empty sum understood when $n=0$.  This
formula, not an analogy with an orbit, is the primitive F-series of the later
papers.

\subsection{The base-\texorpdfstring{$p$}{p} digit compactum}

For any integer $p\geq2$, put
\begin{equation}
\label{eq:composite-p-adics}
 \Z_p^{\mathrm{dig}}=\varprojlim_N\Z/p^N\Z.
\end{equation}
Every element has a unique digit expansion
$\mathfrak z=\sum_{a\geq0}d_a(\mathfrak z)p^a$ with
$d_a(\mathfrak z)\in\mathcal A_p$.  If $p$ is prime this is the usual ring
$\Z_p$.  For composite $p$, \cref{eq:composite-p-adics} is still the compact
base-$p$ digit ring, but it is not a discrete valuation ring and should not
be treated as a field's ring of integers.

\begin{proposition}[digits, the inverse limit, and Haar probability]
\label{prop:digit-compactum-haar}
Give $\mathcal A_p^{\N_0}$ the product topology, with each copy of
$\mathcal A_p$ discrete.  The map
\[
 \Phi_p:\mathcal A_p^{\N_0}\longrightarrow\Z_p^{\mathrm{dig}},
 \qquad
 (d_a)_{a\geq0}\longmapsto
 \left(\sum_{a=0}^{N-1}d_ap^a\bmod p^N\right)_{N\geq1}
\]
is a homeomorphism.  Under $\Phi_p$, normalized Haar probability $m_p$ on
$\Z_p^{\mathrm{dig}}$ is the product of the uniform probability measures on
$\mathcal A_p$.  Equivalently, for every $N\geq1$ and
$e_0,\ldots,e_{N-1}\in\mathcal A_p$,
\begin{equation}
\label{eq:digit-cylinder-mass}
 m_p\{\mathfrak z:d_0(\mathfrak z)=e_0,\ldots,
                   d_{N-1}(\mathfrak z)=e_{N-1}\}=p^{-N}.
\end{equation}
Thus the coordinate digits are independent and each is uniform on
$\mathcal A_p$.
\end{proposition}

\begin{proof}
The displayed residues are compatible under reduction, so $\Phi_p$ is
well-defined.  Conversely, take a compatible sequence
$(x_N\bmod p^N)_N$.  Its reduction modulo $p$ determines a unique
$d_0\in\mathcal A_p$.  After $d_0,\ldots,d_{N-1}$ have been determined,
compatibility makes
\[
 x_{N+1}-\sum_{a=0}^{N-1}d_ap^a
\]
divisible by $p^N$, and its quotient modulo $p$ determines the unique next
digit $d_N$.  This constructs the inverse of $\Phi_p$.  Both maps are
continuous because each output coordinate or digit depends on only finitely
many input coordinates; alternatively, the continuous bijection is a
homeomorphism because the product is compact and the inverse limit is
Hausdorff.

Fixing the first $N$ digits specifies one coset of the open subgroup
$p^N\Z_p^{\mathrm{dig}}$.  That subgroup has index $p^N$, so translation
invariance and total mass one give every such cylinder mass $p^{-N}$.
These are exactly the finite-dimensional cylinder probabilities of the
uniform product measure, proving the measure identity and independence.
\end{proof}

Write
\begin{align}
 [\mathfrak z]_{p^N}&=\sum_{a=0}^{N-1}d_a(\mathfrak z)p^a,\\
 \sh(\mathfrak z)&=\frac{\mathfrak z-d_0(\mathfrak z)}p
                  =\sum_{a\geq0}d_{a+1}(\mathfrak z)p^a.
\end{align}
The first is the canonical truncation and the second is the left digit shift.
The digit and measure assertions used below are therefore consequences of
\cref{prop:digit-compactum-haar}, including when $p$ is composite.

\subsection{The primitive F-series}

Assume throughout this subsection that $H$ is centered.  Fix a place $\ell$
of $K$ and view all $r_j$ as continuous operators on the completion
$K_\ell$.  Let
\begin{equation}
\label{eq:partial-F-series}
 S_N(\mathfrak z)=X_H([\mathfrak z]_{p^N})
 =\sum_{a=0}^{N-1}
   \left(\prod_{k=0}^{a-1}r_{d_k(\mathfrak z)}\right)
   c_{d_a(\mathfrak z)}.
\end{equation}
Here and below the empty product is $I$.  When $S_N(\mathfrak z)$ converges
in $K_\ell$, define
\begin{equation}
\label{eq:rising-extension}
 X_{H,\ell}(\mathfrak z)=\lim_{N\to\infty}S_N(\mathfrak z).
\end{equation}
This is a \emph{rising extension}: its values are prescribed by limits along
the canonical truncation sequence.  It is not a topology-free interpolation.

\begin{proposition}[pointwise convergence criterion]
\label{prop:numen-pointwise-convergence}
Let $C=\max_j|c_j|_\ell$.  Suppose that for a given $\mathfrak z$ there are
$\eta>0$ and $A\in\R$ such that
\begin{equation}
\label{eq:negative-lyapunov}
 \sum_{k=0}^{a-1}\log\|r_{d_k(\mathfrak z)}\|_\ell
 \leq A-\eta a
\end{equation}
for all sufficiently large $a$.  Then the series in
\cref{eq:partial-F-series} converges absolutely in an archimedean completion
and converges in a nonarchimedean completion.  The same estimate makes the
remainder term
$\bigl(\prod_{k<N}r_{d_k}\bigr)a_0$ tend to zero for every fixed
$a_0\in K_\ell$.
\end{proposition}

\begin{proof}
The norm of the $a$th summand is at most
$C\exp(A-\eta a)$.  This is summable in the archimedean case and tends to
zero in the nonarchimedean case; a series in a complete nonarchimedean field
converges exactly when its terms tend to zero.  The identical product bound
applies to the remainder.
\end{proof}

\begin{sourceissue}[the pointwise digit-density criterion]
For a fixed digit sequence, write
\[
 \nu_{j,N}=\frac1N\#\{0\leq k<N:d_k(\mathfrak z)=j\},
 \qquad \alpha_j=\log\|r_j\|_\ell.
\]
The upper Lyapunov exponent is controlled by
\begin{equation}
\label{eq:density-lyapunov-upper-bound}
 \limsup_{N\to\infty}\sum_j\alpha_j\nu_{j,N}
 \leq
 \sum_{\alpha_j<0}\alpha_j\liminf_{N\to\infty}\nu_{j,N}
 +\sum_{\alpha_j>0}\alpha_j\limsup_{N\to\infty}\nu_{j,N}.
\end{equation}
For a negative coefficient, an upper bound on $\alpha_j\nu_{j,N}$ requires a
lower bound on the frequency, hence a liminf; multiplication by
$\alpha_j<0$ reverses the inequality.  For a positive coefficient, it
requires an upper bound on the frequency, hence a limsup.  Terms with
$\alpha_j=0$ are immaterial.  The pointwise criterion in the 2026 source
reverses these choices, using a limsup for contractive branches and a liminf
for expansive branches.  Its displayed quantity is therefore a lower, not
an upper, control and its negativity is not a sufficient convergence
condition.  We avoid that error here by assuming the direct eventual bound
\cref{eq:negative-lyapunov}; negativity of the right side of
\cref{eq:density-lyapunov-upper-bound} is one valid way to obtain that bound.
\end{sourceissue}

\begin{theorem}[almost-everywhere extension]
\label{thm:numen-ae-extension}
Assume $H$ is centered and proper and
\begin{equation}
\label{eq:geometric-mean-contraction}
 \prod_{j=0}^{p-1}\|r_j\|_\ell<1.
\end{equation}
Then \cref{eq:rising-extension} converges for Haar-almost every
$\mathfrak z\in\Z_p^{\mathrm{dig}}$ and defines a measurable
$K_\ell$-valued function on its conull domain.  On that domain it obeys
\begin{equation}
\label{eq:numen-FE-adic}
 X_{H,\ell}(p\mathfrak z+j)
 =r_jX_{H,\ell}(\mathfrak z)+c_j.
\end{equation}
\end{theorem}

\begin{proof}
For Haar-almost every digit sequence, the strong law of large numbers gives
\[
 \frac1N\sum_{k=0}^{N-1}\log\|r_{d_k(\mathfrak z)}\|_\ell
 \longrightarrow \frac1p\sum_{j=0}^{p-1}\log\|r_j\|_\ell<0.
\]
Thus \cref{eq:negative-lyapunov} holds eventually, and
\cref{prop:numen-pointwise-convergence} applies.  Each $S_N$ is locally
constant and therefore measurable; its pointwise limit is measurable.
Finally, the digits of $p\mathfrak z+j$ are $j,d_0(\mathfrak z),\ldots$.
Separate the first summand in \cref{eq:partial-F-series} and pass to the
limit to obtain \cref{eq:numen-FE-adic}.
\end{proof}

\begin{theorem}[uniform contraction]
\label{thm:numen-uniform-contraction}
Suppose
\[
 \rho=\max_{0\leq j<p}\|r_j\|_\ell<1.
\]
Then the series in \cref{eq:partial-F-series} converges uniformly on
$\Z_p^{\mathrm{dig}}$.  Its sum is continuous and satisfies
\cref{eq:numen-FE-adic}.  More precisely, in the archimedean case,
\begin{equation}
\label{eq:uniform-tail-bound}
 \sup_{\mathfrak z}|X_{H,\ell}(\mathfrak z)-S_N(\mathfrak z)|_\ell
 \leq \frac{C\rho^N}{1-\rho},
\end{equation}
while in the nonarchimedean case the right side may be replaced by
$C\rho^N$.
\end{theorem}

\begin{proof}
The $a$th summand has norm at most $C\rho^a$, uniformly in
$\mathfrak z$.  The ordinary and ultrametric tail estimates give
\cref{eq:uniform-tail-bound}.  Uniform limits of the locally constant
functions $S_N$ are continuous.  The functional equation follows exactly as
in \cref{thm:numen-ae-extension}.
\end{proof}

\begin{sourceissue}[centeredness and uniqueness in the extension lemma]
The 2026 extension lemma assumes properness but not centeredness, while its
displayed F-series is the translation term $H_{\mathbf j}(0)$ and therefore
uses terminal value $0$.  For a noncentered proper Hydra, the unique solution
on $\N_0$ instead has terminal value
$a=(I-r_0)^{-1}c_0$ and finite values
$M_H(\mathbf d)a+X_H(\mathbf d)$ as in
\cref{eq:finite-general-solution}.  The extra product term disappears under
a contraction hypothesis, but it cannot be omitted formally.  Likewise, a
functional equation on all infinite digit strings need not determine an
arbitrary function without a rising-limit, regularity, or boundedness
condition.  The uniqueness asserted here is uniqueness of the prescribed
rising extension, and under uniform contraction uniqueness in the natural
continuous bounded class.
\end{sourceissue}

\subsection{The shortened \texorpdfstring{$qx+1$}{qx+1} numen}

For the Hydra $T_q$, write $\chi_q=X_{T_q}$.  From
\cref{eq:Tq-branches,eq:numen-FE-N},
\begin{align}
 \chi_q(2n)&=\frac12\chi_q(n),\label{eq:chi-even}\\
 \chi_q(2n+1)&=\frac{q\chi_q(n)+1}{2},\label{eq:chi-odd}
\end{align}
with $\chi_q(0)=0$.  If
$\operatorname{dig}_2(n)=(d_0,\ldots,d_{L-1})$ and
$s_a=d_0+\cdots+d_{a-1}$, then
\begin{align}
 M_q(n)&=\frac{q^{d_0+\cdots+d_{L-1}}}{2^L},
 \label{eq:Mq-explicit}\\
 \chi_q(n)&=\sum_{\substack{0\leq a<L\\d_a=1}}
               \frac{q^{s_a}}{2^{a+1}}.
 \label{eq:chi-explicit}
\end{align}
Here $L$ is the length of the canonical shortest binary word; in particular,
$M_q(0)=1$.  Appending zeroes would multiply the word multiplier by extra
factors of $1/2$, so a padded representative must not be used in
\cref{eq:Mq-explicit}.
For an odd prime $q$, the same series converges in $\Q_q$ for every
$\mathfrak z\in\Z_2$: it is finite when the binary expansion has finitely
many ones, and otherwise its nonzero terms have $q$-adic norms
$q^{-s_a}\to0$.  Thus
\begin{equation}
\label{eq:chi-q-adic-series}
 \chi_q(\mathfrak z)=
 \sum_{\substack{a\geq0\\d_a(\mathfrak z)=1}}
 \frac{q^{s_a(\mathfrak z)}}{2^{a+1}}
 \in\Z_q.
\end{equation}
It is rising-continuous everywhere.  It is $2$-adic-to-$q$-adic continuous
at every binary sequence with infinitely many ones, but not at the
nonnegative integers: if $n\in\N_0$ has $k=\#_1(n)$ ones, a new $1$ placed
arbitrarily far beyond its last digit contributes a term of the fixed
nonzero $q$-adic norm $q^{-k}$, independent of that remote position (it is a
unit only when $n=0$).  This pointwise distinction is the prototype for the
later theory of variable topologies of convergence.
